% Options for packages loaded elsewhere
\PassOptionsToPackage{unicode}{hyperref}
\PassOptionsToPackage{hyphens}{url}
\documentclass[
]{article}
\usepackage{xcolor}
\usepackage[margin=1in]{geometry}
\usepackage{amsmath,amssymb}
\setcounter{secnumdepth}{-\maxdimen} % remove section numbering
\usepackage{iftex}
\ifPDFTeX
  \usepackage[T1]{fontenc}
  \usepackage[utf8]{inputenc}
  \usepackage{textcomp} % provide euro and other symbols
\else % if luatex or xetex
  \usepackage{unicode-math} % this also loads fontspec
  \defaultfontfeatures{Scale=MatchLowercase}
  \defaultfontfeatures[\rmfamily]{Ligatures=TeX,Scale=1}
\fi
\usepackage{lmodern}
\ifPDFTeX\else
  % xetex/luatex font selection
\fi
% Use upquote if available, for straight quotes in verbatim environments
\IfFileExists{upquote.sty}{\usepackage{upquote}}{}
\IfFileExists{microtype.sty}{% use microtype if available
  \usepackage[]{microtype}
  \UseMicrotypeSet[protrusion]{basicmath} % disable protrusion for tt fonts
}{}
\makeatletter
\@ifundefined{KOMAClassName}{% if non-KOMA class
  \IfFileExists{parskip.sty}{%
    \usepackage{parskip}
  }{% else
    \setlength{\parindent}{0pt}
    \setlength{\parskip}{6pt plus 2pt minus 1pt}}
}{% if KOMA class
  \KOMAoptions{parskip=half}}
\makeatother
\usepackage{color}
\usepackage{fancyvrb}
\newcommand{\VerbBar}{|}
\newcommand{\VERB}{\Verb[commandchars=\\\{\}]}
\DefineVerbatimEnvironment{Highlighting}{Verbatim}{commandchars=\\\{\}}
% Add ',fontsize=\small' for more characters per line
\usepackage{framed}
\definecolor{shadecolor}{RGB}{248,248,248}
\newenvironment{Shaded}{\begin{snugshade}}{\end{snugshade}}
\newcommand{\AlertTok}[1]{\textcolor[rgb]{0.94,0.16,0.16}{#1}}
\newcommand{\AnnotationTok}[1]{\textcolor[rgb]{0.56,0.35,0.01}{\textbf{\textit{#1}}}}
\newcommand{\AttributeTok}[1]{\textcolor[rgb]{0.13,0.29,0.53}{#1}}
\newcommand{\BaseNTok}[1]{\textcolor[rgb]{0.00,0.00,0.81}{#1}}
\newcommand{\BuiltInTok}[1]{#1}
\newcommand{\CharTok}[1]{\textcolor[rgb]{0.31,0.60,0.02}{#1}}
\newcommand{\CommentTok}[1]{\textcolor[rgb]{0.56,0.35,0.01}{\textit{#1}}}
\newcommand{\CommentVarTok}[1]{\textcolor[rgb]{0.56,0.35,0.01}{\textbf{\textit{#1}}}}
\newcommand{\ConstantTok}[1]{\textcolor[rgb]{0.56,0.35,0.01}{#1}}
\newcommand{\ControlFlowTok}[1]{\textcolor[rgb]{0.13,0.29,0.53}{\textbf{#1}}}
\newcommand{\DataTypeTok}[1]{\textcolor[rgb]{0.13,0.29,0.53}{#1}}
\newcommand{\DecValTok}[1]{\textcolor[rgb]{0.00,0.00,0.81}{#1}}
\newcommand{\DocumentationTok}[1]{\textcolor[rgb]{0.56,0.35,0.01}{\textbf{\textit{#1}}}}
\newcommand{\ErrorTok}[1]{\textcolor[rgb]{0.64,0.00,0.00}{\textbf{#1}}}
\newcommand{\ExtensionTok}[1]{#1}
\newcommand{\FloatTok}[1]{\textcolor[rgb]{0.00,0.00,0.81}{#1}}
\newcommand{\FunctionTok}[1]{\textcolor[rgb]{0.13,0.29,0.53}{\textbf{#1}}}
\newcommand{\ImportTok}[1]{#1}
\newcommand{\InformationTok}[1]{\textcolor[rgb]{0.56,0.35,0.01}{\textbf{\textit{#1}}}}
\newcommand{\KeywordTok}[1]{\textcolor[rgb]{0.13,0.29,0.53}{\textbf{#1}}}
\newcommand{\NormalTok}[1]{#1}
\newcommand{\OperatorTok}[1]{\textcolor[rgb]{0.81,0.36,0.00}{\textbf{#1}}}
\newcommand{\OtherTok}[1]{\textcolor[rgb]{0.56,0.35,0.01}{#1}}
\newcommand{\PreprocessorTok}[1]{\textcolor[rgb]{0.56,0.35,0.01}{\textit{#1}}}
\newcommand{\RegionMarkerTok}[1]{#1}
\newcommand{\SpecialCharTok}[1]{\textcolor[rgb]{0.81,0.36,0.00}{\textbf{#1}}}
\newcommand{\SpecialStringTok}[1]{\textcolor[rgb]{0.31,0.60,0.02}{#1}}
\newcommand{\StringTok}[1]{\textcolor[rgb]{0.31,0.60,0.02}{#1}}
\newcommand{\VariableTok}[1]{\textcolor[rgb]{0.00,0.00,0.00}{#1}}
\newcommand{\VerbatimStringTok}[1]{\textcolor[rgb]{0.31,0.60,0.02}{#1}}
\newcommand{\WarningTok}[1]{\textcolor[rgb]{0.56,0.35,0.01}{\textbf{\textit{#1}}}}
\usepackage{graphicx}
\makeatletter
\newsavebox\pandoc@box
\newcommand*\pandocbounded[1]{% scales image to fit in text height/width
  \sbox\pandoc@box{#1}%
  \Gscale@div\@tempa{\textheight}{\dimexpr\ht\pandoc@box+\dp\pandoc@box\relax}%
  \Gscale@div\@tempb{\linewidth}{\wd\pandoc@box}%
  \ifdim\@tempb\p@<\@tempa\p@\let\@tempa\@tempb\fi% select the smaller of both
  \ifdim\@tempa\p@<\p@\scalebox{\@tempa}{\usebox\pandoc@box}%
  \else\usebox{\pandoc@box}%
  \fi%
}
% Set default figure placement to htbp
\def\fps@figure{htbp}
\makeatother
\setlength{\emergencystretch}{3em} % prevent overfull lines
\providecommand{\tightlist}{%
  \setlength{\itemsep}{0pt}\setlength{\parskip}{0pt}}
\usepackage{bookmark}
\IfFileExists{xurl.sty}{\usepackage{xurl}}{} % add URL line breaks if available
\urlstyle{same}
\hypersetup{
  pdftitle={The Gamma-Poisson Model},
  hidelinks,
  pdfcreator={LaTeX via pandoc}}

\title{The Gamma-Poisson Model}
\author{}
\date{\vspace{-2.5em}2026-04-17}

\begin{document}
\maketitle

\subsection{Introduction}\label{introduction}

The Gamma--Poisson conjugate pair is the count-data analogue of the
Beta--Binomial model: a Gamma prior on a Poisson rate yields a Gamma
posterior whose parameters are obtained by adding the observed event
count and the observed exposure to the prior shape and rate
respectively. It is the natural starting point whenever you are
inferring an event rate per unit time, area, or population, and it
generalises directly to Poisson regression once predictors enter the
linear predictor.

\subsection{Prerequisites}\label{prerequisites}

A working understanding of the Poisson distribution, the Gamma
distribution, and the concept of conjugacy.

\subsection{Theory}\label{theory}

With prior \(\lambda \sim \mathrm{Gamma}(\alpha_0, \beta_0)\) and
likelihood \(y_i \sim \mathrm{Poisson}(\lambda t_i)\) where \(t_i\) is
the exposure for observation \(i\), the posterior is

\[\lambda \mid y \sim \mathrm{Gamma}\!\left(\alpha_0 + \sum_{i} y_i, \, \beta_0 + \sum_{i} t_i\right).\]

The prior shape \(\alpha_0\) behaves like a pseudo-count of events and
the prior rate \(\beta_0\) like a pseudo-exposure, so a
\(\mathrm{Gamma}(2, 1)\) prior is equivalent to two prior events
observed in one unit of exposure. The posterior mean is
\((\alpha_0 + \sum y_i) / (\beta_0 + \sum t_i)\) and the posterior
variance is the posterior mean divided by the posterior rate. The
Poisson--Gamma compound distribution is a Negative-Binomial, which gives
the posterior predictive distribution in closed form.

\subsection{Assumptions}\label{assumptions}

Independent Poisson counts conditional on the rate, an exposure \(t_i\)
that is known and not estimated from the same data, and a constant rate
over each observation window (after conditioning on any covariates).
Overdispersion relative to Poisson --- variance exceeding the mean ---
motivates a Negative-Binomial likelihood instead.

\subsection{R Implementation}\label{r-implementation}

\begin{Shaded}
\begin{Highlighting}[]
\CommentTok{\# Prior Gamma(2, 1) {-}{-} 2 events per 1 unit of exposure}
\NormalTok{alpha0 }\OtherTok{\textless{}{-}} \DecValTok{2}\NormalTok{; beta0 }\OtherTok{\textless{}{-}} \DecValTok{1}

\CommentTok{\# Data: 15 events in 10 units of exposure}
\NormalTok{y }\OtherTok{\textless{}{-}} \DecValTok{15}\NormalTok{; t }\OtherTok{\textless{}{-}} \DecValTok{10}

\NormalTok{alpha\_n }\OtherTok{\textless{}{-}}\NormalTok{ alpha0 }\SpecialCharTok{+}\NormalTok{ y}
\NormalTok{beta\_n  }\OtherTok{\textless{}{-}}\NormalTok{ beta0 }\SpecialCharTok{+}\NormalTok{ t}

\FunctionTok{c}\NormalTok{(}\AttributeTok{posterior\_mean =}\NormalTok{ alpha\_n }\SpecialCharTok{/}\NormalTok{ beta\_n,}
  \AttributeTok{posterior\_sd   =} \FunctionTok{sqrt}\NormalTok{(alpha\_n }\SpecialCharTok{/}\NormalTok{ beta\_n}\SpecialCharTok{\^{}}\DecValTok{2}\NormalTok{),}
  \AttributeTok{ci95 =} \FunctionTok{qgamma}\NormalTok{(}\FunctionTok{c}\NormalTok{(}\FloatTok{0.025}\NormalTok{, }\FloatTok{0.975}\NormalTok{), alpha\_n, beta\_n))}
\end{Highlighting}
\end{Shaded}

\subsection{Output \& Results}\label{output-results}

Closed-form summaries: posterior mean approximately 1.55 events per unit
exposure, posterior SD approximately 0.38, and 95 \% equal-tailed
credible interval roughly \((0.94, 2.28)\). The posterior is
right-skewed because Gamma densities have a heavier right tail than
left, so report the median or HPD when the asymmetry matters.

\subsection{Interpretation}\label{interpretation}

A reporting sentence: ``A \(\mathrm{Gamma}(2, 1)\) prior combined with
15 events in 10 units of exposure yielded a \(\mathrm{Gamma}(17, 11)\)
posterior with mean 1.55 events per unit exposure (95 \% credible
interval 0.94 to 2.28); the prior contributed the equivalent of two
pseudo-events and one pseudo-unit of exposure, which the data
overwhelmed.'' Quote both the posterior mean and the credible interval,
and note any skew.

\subsection{Practical Tips}\label{practical-tips}

\begin{itemize}
\tightlist
\item
  Weakly informative defaults: \(\mathrm{Gamma}(0.5, 0.5)\) or
  \(\mathrm{Gamma}(1, 1)\); the latter has prior mean 1 and prior
  variance 1.
\item
  For Poisson regression, place priors on regression coefficients on the
  log-rate scale (typically \(\mathcal N(0, 1)\) or
  \(\mathcal N(0, 2.5)\)) rather than on the rate itself; the linear
  predictor is more naturally constrained.
\item
  The posterior predictive distribution is a Negative-Binomial:
  \(\tilde y \mid y \sim \mathrm{NegBinom}\bigl(\alpha_n, t_{\text{new}} / (\beta_n + t_{\text{new}})\bigr)\),
  so prediction intervals come for free without simulation.
\item
  Persistent overdispersion after fitting motivates a Negative-Binomial
  likelihood with a separate dispersion parameter; the Gamma--Poisson is
  its conjugate root.
\item
  For zero-inflation (more zeros than the Poisson predicts), augment
  with a zero-inflation component rather than tightening the prior.
\end{itemize}

\subsection{R Packages Used}\label{r-packages-used}

Base R suffices for the conjugate update (\texttt{dgamma},
\texttt{qgamma}); \texttt{bayestestR} for tidy summaries; \texttt{brms}
with \texttt{family\ =\ poisson()} (or \texttt{negbinomial()}) when
covariate-dependent rates are needed; \texttt{rstanarm::stan\_glm} with
\texttt{family\ =\ poisson} for a quick GLM-style alternative.

\subsection{For Reviewers}\label{for-reviewers}

\textbf{What to look for in a paper using this method.}

\begin{itemize}
\tightlist
\item
  \textbf{Common misapplications.}
\item
  \textbf{Diagnostics that should be reported but often aren't.}
\item
  \textbf{Red flags in tables and figures.}
\item
  \textbf{What to verify.}
\item
  \textbf{What an adequate Methods paragraph must contain.}
\end{itemize}

\subsection{See also --- labs in this
chapter}\label{see-also-labs-in-this-chapter}

\begin{itemize}
\tightlist
\item
  \href{labs/lab-bayesian-thinking-control-vs-decision-error.Rmd}{Bayesian
  thinking; α control vs decision error}
\item
  \href{labs/lab-brms-stan-loo-hierarchical-models.Rmd}{brms / Stan,
  LOO, hierarchical models}
\end{itemize}

\end{document}
